\documentclass[12pt,oneside]{book}

% Basic packages
\usepackage[utf8]{inputenc}
\usepackage[T1]{fontenc}
\usepackage{lmodern}
\usepackage[margin=1in]{geometry}
\usepackage{setspace}
\usepackage{amsmath,amssymb}
\usepackage{booktabs}
\usepackage{array}
\usepackage{hyperref}
\usepackage{xcolor}
\usepackage{enumitem}

\hypersetup{
    colorlinks=true,
    linkcolor=blue,
    urlcolor=blue,
    citecolor=blue
}

\onehalfspacing

% Simple table column type for left-aligned text with fixed width
\newcolumntype{L}[1]{>{\raggedright\arraybackslash}p{#1}}

\begin{document}

\frontmatter

\title{Comprehensive High School Course Syllabi\\
\large Complete Edition}
\author{Institutional-Grade Curriculum Documentation}
\date{January 2026}

\maketitle

\thispagestyle{empty}
\vspace*{\fill}
\begin{center}
\textbf{Complete Curriculum Guide for All High School Courses}\\[1em]
Including AP, Honors, and Dual Credit Programs
\end{center}

\vspace{1em}
\noindent Includes:
\begin{itemize}[leftmargin=2em]
    \item Core Courses: English, Math, Science, Social Studies
    \item Advanced Placement (AP) Courses
    \item Honors Courses
    \item Dual Credit Courses
    \item Mathematics: Pre-Calculus, Calculus I, Statistics
    \item Complete AP Course Syllabi
    \item College-Level Course Descriptions
\end{itemize}

\vspace{1em}
\noindent Developed in accordance with:
\begin{itemize}[leftmargin=2em]
    \item Common Core State Standards
    \item Next Generation Science Standards (NGSS)
    \item College Board AP Frameworks
    \item International Baccalaureate Requirements
    \item State Dual Credit Articulation Agreements
\end{itemize}
\vspace*{\fill}

\tableofcontents

\chapter*{Preface}
\addcontentsline{toc}{chapter}{Preface}

This comprehensive curriculum documentation represents institutional-grade syllabi for all major
high school courses including regular, honors, Advanced Placement (AP), and dual credit options.
Each syllabus has been developed to meet national and state standards while providing flexibility
for local adaptation.

\section*{Purpose}
This document serves multiple audiences:
\begin{itemize}[leftmargin=2em]
    \item \textbf{Students}: Preview course content and understand expectations at all levels
    \item \textbf{Parents}: Support student learning and make informed course selections
    \item \textbf{Teachers}: Comprehensive curriculum framework for instruction across all levels
    \item \textbf{Administrators}: Documentation for accreditation, program evaluation, and college partnerships
    \item \textbf{Counselors}: Guide students in appropriate course placement and college preparation
\end{itemize}

\section*{Course Level Distinctions}
\subsection*{Regular Courses}
Standard high school courses meeting state graduation requirements with grade-appropriate rigor.

\subsection*{Honors Courses}
Accelerated courses with:
\begin{itemize}[leftmargin=2em]
    \item Faster pacing and greater depth
    \item More independent work and critical thinking
    \item Advanced reading and writing expectations
    \item Weighted GPA (typically +0.5 points)
    \item Preparation for AP-level coursework
\end{itemize}

\subsection*{Advanced Placement (AP) Courses}
College-level courses with:
\begin{itemize}[leftmargin=2em]
    \item Curriculum designed by College Board
    \item Opportunity to earn college credit via AP exam
    \item Weighted GPA (typically +1.0 point)
    \item Rigorous content and high expectations
    \item Exams administered in May each year
\end{itemize}

\subsection*{Dual Credit Courses}
Courses offering both high school and college credit:
\begin{itemize}[leftmargin=2em]
    \item Taught by qualified instructors (often college-approved)
    \item Follow college curriculum and standards
    \item Transcripted by partner college/university
    \item May have additional fees or requirements
    \item Credits transferable to many institutions
\end{itemize}

\section*{Standards Alignment}
All syllabi align with:
\begin{itemize}[leftmargin=2em]
    \item Common Core State Standards (where applicable)
    \item Next Generation Science Standards (NGSS)
    \item National Council for the Social Studies (NCSS) standards
    \item College Board Advanced Placement frameworks
    \item International Baccalaureate Diploma Programme requirements
    \item State dual credit articulation agreements
    \item Partner college/university course competencies
\end{itemize}

\mainmatter

% =========================
% Part I: English/Language Arts
% =========================
\part{English/Language Arts}

\chapter{English 9 (Freshman English)}

\noindent\textbf{Course:} English 9\\
\textbf{Grade Level:} 9th Grade\\
\textbf{Duration:} Full Year (36 weeks)\\
\textbf{Credits:} 1.0 Credit\\
\textbf{Prerequisites:} None

\section{Course Description}
English 9 introduces students to high school-level literary analysis, writing, and communication
skills through the study of fiction, nonfiction, poetry, and drama. Students engage with diverse
texts organized thematically while developing critical thinking and analytical writing abilities.

\section{Course Goals}
Upon completion of this course, students will be able to:
\begin{itemize}[leftmargin=2em]
    \item Develop literary analysis skills through close reading of complex texts
    \item Strengthen expository and creative writing abilities
    \item Build vocabulary and grammar proficiency
    \item Practice effective oral communication and presentation skills
    \item Engage in Socratic seminars and collaborative discussions
    \item Apply literary devices and elements to text analysis
    \item Write thesis-driven literary analysis essays
\end{itemize}

\section{Grading Breakdown}
\begin{center}
\begin{tabular}{L{0.55\linewidth}L{0.25\linewidth}}
\toprule
\textbf{Category} & \textbf{Percentage} \\
\midrule
Tests and Quizzes & 30\% \\
Essays and Writing Assignments & 30\% \\
Projects and Presentations & 20\% \\
Class Participation and Discussion & 10\% \\
Homework and Daily Work & 10\% \\
\bottomrule
\end{tabular}
\end{center}

\chapter{Honors English 9}

\noindent\textbf{HONORS COURSE:} Honors English 9\\
\textbf{Grade Level:} 9th Grade\\
\textbf{Duration:} Full Year (36 weeks)\\
\textbf{Credits:} 1.0 Credit (Weighted 4.5 scale)\\
\textbf{Prerequisites:} Teacher recommendation and/or 8th grade English grade of B+ or higher

\section{Course Description}
Honors English 9 provides an accelerated and in-depth exploration of literary analysis, rhetoric,
and composition. Students read more complex texts, write more sophisticated essays, and engage
in higher-level critical thinking than in regular English 9.

\section{Distinctions from Regular English 9}
\begin{itemize}[leftmargin=2em]
    \item \textbf{Pace:} Faster progression through units with less scaffolding
    \item \textbf{Reading:} Additional novels, more challenging texts, independent reading requirements
    \item \textbf{Writing:} Longer, more sophisticated essays with higher expectations for style and analysis
    \item \textbf{Discussion:} More frequent Socratic seminars and student-led discussions
    \item \textbf{Independence:} Greater emphasis on self-directed learning and research
\end{itemize}

\section{Additional Texts (Beyond English 9 Core)}
\begin{itemize}[leftmargin=2em]
    \item \emph{Night} by Elie Wiesel
    \item \emph{Fahrenheit 451} by Ray Bradbury
    \item Selection of more challenging poetry and short stories
    \item Independent reading novels (2--3 additional books)
\end{itemize}

\section{Grading Breakdown}
\begin{center}
\begin{tabular}{L{0.55\linewidth}L{0.25\linewidth}}
\toprule
\textbf{Category} & \textbf{Percentage} \\
\midrule
Tests and Essays & 35\% \\
Writing Assignments and Projects & 30\% \\
Independent Reading and Analysis & 15\% \\
Class Participation and Discussion & 10\% \\
Homework and Daily Work & 10\% \\
\bottomrule
\end{tabular}
\end{center}

\chapter{AP English Language and Composition}

\noindent\textbf{AP COURSE:} AP English Language and Composition\\
\textbf{Grade Level:} 11th Grade\\
\textbf{Duration:} Full Year (36 weeks)\\
\textbf{Credits:} 1.0 Credit (Weighted 5.0 scale)\\
\textbf{Prerequisites:} Honors English 10 or English 10 with teacher recommendation

\section{Course Description}
AP English Language and Composition focuses on rhetorical analysis, argumentation, and synthesis.
Students analyze nonfiction texts from various time periods and disciplines while developing sophisticated
writing skills across multiple modes. This course prepares students for the AP Language and Composition
exam in May.

\section{Course Goals}
Students will:
\begin{itemize}[leftmargin=2em]
    \item Analyze rhetorical strategies in nonfiction texts
    \item Construct sophisticated arguments using evidence
    \item Write in multiple modes: rhetorical analysis, argument, synthesis
    \item Understand classical and contemporary rhetoric
    \item Develop a mature writing style
    \item Prepare for AP exam and college-level writing
    \item Engage critically with complex ideas and perspectives
\end{itemize}

\section{AP Exam Components}
\subsection*{Multiple Choice (45\% of score)}
\begin{itemize}[leftmargin=2em]
    \item 45 questions in 60 minutes
    \item Reading passages with rhetorical analysis questions
\end{itemize}

\subsection*{Free Response (55\% of score)}
\begin{itemize}[leftmargin=2em]
    \item \textbf{Synthesis Essay (15\%):} Develop a position using provided sources
    \item \textbf{Rhetorical Analysis Essay (15\%):} Analyze an author's rhetorical choices
    \item \textbf{Argument Essay (25\%):} Construct an evidence-based argument on a given topic
    \item Total time: 2 hours 15 minutes
\end{itemize}

\section{Units of Study}
\subsection*{Unit 1: Introduction to Rhetoric and Rhetorical Analysis (4 weeks)}
Topics:
\begin{itemize}[leftmargin=2em]
    \item Classical rhetoric: Aristotle's \emph{Rhetoric}
    \item Rhetorical triangle: speaker, audience, subject
    \item Rhetorical appeals: ethos, pathos, logos
    \item Rhetorical situation: exigence, context, constraints
    \item SOAPS (Speaker, Occasion, Audience, Purpose, Subject)
    \item Rhetorical devices and strategies
    \item Close reading of nonfiction
\end{itemize}

\subsection*{Unit 2: Analysis of Argument (5 weeks)}
Topics:
\begin{itemize}[leftmargin=2em]
    \item Components of argument: claim, evidence, warrant, backing, qualifier, rebuttal
    \item Toulmin model of argument
    \item Logical fallacies
    \item Evaluating evidence
    \item Analyzing visual arguments
    \item Analyzing speeches and essays
\end{itemize}

\subsection*{Unit 3: Building Arguments (5 weeks)}
Topics:
\begin{itemize}[leftmargin=2em]
    \item Crafting thesis statements
    \item Using evidence effectively
    \item Addressing counterarguments
    \item Qualifying claims
    \item Organizing arguments
    \item Developing voice and style
\end{itemize}

\subsection*{Unit 4: Synthesis and Research (5 weeks)}
Topics:
\begin{itemize}[leftmargin=2em]
    \item Synthesizing multiple sources
    \item Evaluating source credibility
    \item Integrating sources smoothly
    \item MLA documentation
    \item Avoiding plagiarism
    \item Synthesis essay practice
\end{itemize}

\subsection*{Unit 5: American Voices and Perspectives (6 weeks)}
Topics and readings:
\begin{itemize}[leftmargin=2em]
    \item Colonial and Revolutionary rhetoric
    \item 19th century reform movements
    \item Civil Rights rhetoric
    \item Contemporary social commentary
    \item Diverse voices and perspectives
\end{itemize}

\subsection*{Unit 6: Satire, Humor, and Irony (3 weeks)}
\subsection*{Unit 7: Science and Nature Writing (3 weeks)}
\subsection*{Unit 8: AP Exam Preparation and Timed Writing (5 weeks)}

\section{Major Writing Assignments}
Students complete approximately:
\begin{itemize}[leftmargin=2em]
    \item 8--10 full rhetorical analysis essays
    \item 8--10 argument essays
    \item 6--8 synthesis essays
    \item Weekly timed writing practice (especially second semester)
    \item Personal essays and reflections
    \item Research project with synthesis
\end{itemize}

\section{Grading Breakdown}
\begin{center}
\begin{tabular}{L{0.55\linewidth}L{0.25\linewidth}}
\toprule
\textbf{Category} & \textbf{Percentage} \\
\midrule
Essays and Timed Writes & 40\% \\
Tests and Quizzes & 25\% \\
Rhetorical Analysis Assignments & 20\% \\
Class Participation and Discussion & 10\% \\
Homework and Reading Checks & 5\% \\
\bottomrule
\end{tabular}
\end{center}

\section{Summer Assignment}
Students must complete summer reading and a rhetorical analysis assignment before school begins.
Details are provided at course registration.

\chapter{AP English Literature and Composition}

\noindent\textbf{AP COURSE:} AP English Literature and Composition\\
\textbf{Grade Level:} 12th Grade\\
\textbf{Duration:} Full Year (36 weeks)\\
\textbf{Credits:} 1.0 Credit (Weighted 5.0 scale)\\
\textbf{Prerequisites:} AP English Language or Honors English 11 with teacher recommendation

\section{Course Description}
AP English Literature and Composition engages students in careful reading and critical analysis of
literature, including poetry, drama, and prose fiction. Students examine how authors use literary
elements and techniques to create meaning while developing sophisticated interpretive and analytical
writing skills.

\section{AP Exam Components}
\subsection*{Multiple Choice (45\% of score)}
\begin{itemize}[leftmargin=2em]
    \item 55 questions in 60 minutes
    \item Poetry and prose fiction passages with analysis questions
\end{itemize}

\subsection*{Free Response (55\% of score)}
\begin{itemize}[leftmargin=2em]
    \item \textbf{Poetry Analysis:} Analyze a given poem
    \item \textbf{Prose Fiction Analysis:} Analyze a prose passage
    \item \textbf{Literary Argument:} Defend an interpretation using a literary work of choice
\end{itemize}

\section{Units of Study}
\subsection*{Unit 1: Introduction to Literary Analysis (3 weeks)}

\subsection*{Unit 2: Poetry Analysis (6--7 weeks)}
Topics:
\begin{itemize}[leftmargin=2em]
    \item Close reading of poetry
    \item Identifying and analyzing poetic devices
    \item Understanding form and structure
    \item Tone, diction, and imagery
    \item Writing poetry analysis essays
    \item Poets studied: Shakespeare, Donne, Dickinson, Frost, Plath, Hughes, contemporary poets
\end{itemize}

\subsection*{Unit 3: Drama (6--7 weeks)}
Core texts:
\begin{itemize}[leftmargin=2em]
    \item \emph{Hamlet} by William Shakespeare
    \item \emph{A Doll's House} by Henrik Ibsen \textbf{or} \emph{Death of a Salesman} by Arthur Miller
\end{itemize}

\subsection*{Unit 4: 19th Century Novel (5--6 weeks)}
Core text options:
\begin{itemize}[leftmargin=2em]
    \item \emph{Jane Eyre} by Charlotte Brontë
    \item \emph{Wuthering Heights} by Emily Brontë
    \item \emph{Great Expectations} by Charles Dickens
    \item \emph{Pride and Prejudice} by Jane Austen
\end{itemize}

\subsection*{Unit 5: 20th/21st Century Novel (5--6 weeks)}
Core text options:
\begin{itemize}[leftmargin=2em]
    \item \emph{Beloved} by Toni Morrison
    \item \emph{The Great Gatsby} by F. Scott Fitzgerald
    \item \emph{Things Fall Apart} by Chinua Achebe
    \item \emph{One Hundred Years of Solitude} by Gabriel García Márquez
\end{itemize}

\subsection*{Unit 6: Short Fiction and Prose Analysis (4 weeks)}

\subsection*{Unit 7: Open Question Preparation (4 weeks)}
Focus:
\begin{itemize}[leftmargin=2em]
    \item Reading AP-level literature independently
    \item Building a repertoire of literary works
    \item Writing literary arguments
    \item Defending interpretations with textual evidence
\end{itemize}

\subsection*{Unit 8: AP Exam Preparation (5 weeks)}

\section{Grading Breakdown}
\begin{center}
\begin{tabular}{L{0.55\linewidth}L{0.25\linewidth}}
\toprule
\textbf{Category} & \textbf{Percentage} \\
\midrule
Essays and Timed Writes & 40\% \\
Tests and Quizzes & 25\% \\
Literary Analysis Assignments & 20\% \\
Class Participation and Discussion & 10\% \\
Reading Checks and Homework & 5\% \\
\bottomrule
\end{tabular}
\end{center}

% =========================
% Part II: Mathematics
% =========================
\part{Mathematics}

\chapter{Algebra 1}

\noindent\textbf{Course:} Algebra 1\\
\textbf{Grade Level:} 9th--10th Grade\\
\textbf{Duration:} Full Year (36 weeks)\\
\textbf{Credits:} 1.0 Credit\\
\textbf{Prerequisites:} Pre-Algebra or equivalent

\section{Course Description}
Algebra 1 is the foundational high school mathematics course introducing algebraic thinking, functions,
and mathematical modeling.

\section{Grading Breakdown}
\begin{center}
\begin{tabular}{L{0.55\linewidth}L{0.25\linewidth}}
\toprule
\textbf{Category} & \textbf{Percentage} \\
\midrule
Tests & 40\% \\
Quizzes & 25\% \\
Homework & 20\% \\
Projects & 10\% \\
Class Participation & 5\% \\
\bottomrule
\end{tabular}
\end{center}

\chapter{Honors Algebra 1}

\noindent\textbf{HONORS COURSE:} Honors Algebra 1\\
\textbf{Grade Level:} 8th--9th Grade\\
\textbf{Duration:} Full Year (36 weeks)\\
\textbf{Credits:} 1.0 Credit (Weighted 4.5 scale)\\
\textbf{Prerequisites:} Teacher recommendation, B+ or higher in Pre-Algebra

\section{Course Description}
Honors Algebra 1 covers all Algebra 1 content at an accelerated pace with additional enrichment topics,
preparing students for Honors Geometry.

\section{Distinctions from Regular Algebra 1}
\begin{itemize}[leftmargin=2em]
    \item Faster pacing with deeper conceptual understanding
    \item More complex problem-solving and multi-step applications
    \item Introduction to advanced topics (matrices, sequences/series beyond basics)
    \item Competition-level problems (AMC 10 preparation)
    \item More rigorous assessments
    \item Emphasis on mathematical proof and reasoning
\end{itemize}

\section{Enrichment Topics}
\begin{itemize}[leftmargin=2em]
    \item Advanced factoring techniques
    \item Introduction to logarithms
    \item Arithmetic and geometric sequences/series
    \item Basic matrix operations
    \item Advanced word problems and modeling
    \item Competition mathematics problems
\end{itemize}

\chapter{Geometry}

\noindent\textbf{Course:} Geometry\\
\textbf{Grade Level:} 9th--11th Grade\\
\textbf{Duration:} Full Year (36 weeks)\\
\textbf{Credits:} 1.0 Credit\\
\textbf{Prerequisites:} Algebra 1

\section{Course Description}
Geometry explores the properties of shapes, spatial reasoning, and formal mathematical proof.
Students learn to make and prove conjectures about geometric figures while developing logical
thinking and problem-solving skills through constructions, proofs, and applications.

\section{Course Goals}
Students will:
\begin{itemize}[leftmargin=2em]
    \item Understand basic geometric definitions, postulates, and theorems
    \item Construct formal geometric proofs (two-column, paragraph, flow-chart)
    \item Apply geometric concepts to solve real-world problems
    \item Use coordinate geometry to analyze figures algebraically
    \item Understand congruence and similarity transformations
    \item Apply trigonometric ratios to solve problems
    \item Calculate area, surface area, and volume of geometric figures
\end{itemize}

\section{Required Materials}
\begin{itemize}[leftmargin=2em]
    \item Compass and straightedge
    \item Protractor
    \item Ruler
    \item Graphing calculator
    \item Graph paper
    \item Three-ring binder
\end{itemize}

\section{Units of Study}

\subsection*{Unit 1: Basics of Geometry and Construction (3--4 weeks)}
Topics:
\begin{itemize}[leftmargin=2em]
    \item Points, lines, and planes
    \item Measuring and constructing segments
    \item Measuring and constructing angles
    \item Angle relationships (complementary, supplementary, vertical, linear pairs)
    \item Midpoint and distance formulas
    \item Perimeter, circumference, and area (review)
    \item Basic geometric constructions
\end{itemize}

Key constructions:
\begin{itemize}[leftmargin=2em]
    \item Copying a segment
    \item Copying an angle
    \item Bisecting a segment
    \item Bisecting an angle
    \item Constructing perpendicular lines
    \item Constructing parallel lines
\end{itemize}

\subsection*{Unit 2: Logical Reasoning and Proof (3 weeks)}
Topics:
\begin{itemize}[leftmargin=2em]
    \item Conditional statements (if--then)
    \item Converse, inverse, and contrapositive
    \item Deductive reasoning
    \item Biconditional statements
    \item Properties of equality and congruence
    \item Proving statements about segments and angles
    \item Introduction to two-column proofs
\end{itemize}

\subsection*{Unit 3: Parallel and Perpendicular Lines (3--4 weeks)}

\subsection*{Unit 4: Transformations (3--4 weeks)}
Topics:
\begin{itemize}[leftmargin=2em]
    \item Translations: $(x,y)\rightarrow(x+a,y+b)$
    \item Reflections over lines
    \item Rotations about a point
    \item Compositions of transformations
    \item Line and rotational symmetry
    \item Dilations: $(x,y)\rightarrow(kx,ky)$
    \item Rigid motions vs. similarity transformations
\end{itemize}

\subsection*{Unit 5: Congruent Triangles (4--5 weeks)}
Topics:
\begin{itemize}[leftmargin=2em]
    \item Classifying triangles
    \item Triangle angle relationships
    \item Triangle Sum Theorem (angles sum to $180^\circ$)
    \item Congruence postulates: SSS, SAS, ASA, AAS, HL
    \item CPCTC (Corresponding Parts of Congruent Triangles are Congruent)
    \item Isosceles and equilateral triangles
    \item Triangle Inequality Theorem
\end{itemize}

\subsection*{Unit 6: Quadrilaterals and Polygons (4 weeks)}

\subsection*{Unit 7: Similarity (4 weeks)}

\subsection*{Unit 8: Right Triangles and Trigonometry (5--6 weeks)}
Topics:
\begin{itemize}[leftmargin=2em]
    \item Geometric mean
    \item Pythagorean Theorem: $a^2 + b^2 = c^2$
    \item Special right triangles: $45^\circ$--$45^\circ$--$90^\circ$ and $30^\circ$--$60^\circ$--$90^\circ$
    \item Trigonometric ratios:
    \[
    \sin\theta = \frac{\text{opposite}}{\text{hypotenuse}},\quad
    \cos\theta = \frac{\text{adjacent}}{\text{hypotenuse}},\quad
    \tan\theta = \frac{\text{opposite}}{\text{adjacent}}
    \]
    \item Solving right triangles
    \item Angles of elevation and depression
    \item Law of Sines and Law of Cosines (if time permits)
\end{itemize}

\subsection*{Unit 9: Circles (4--5 weeks)}
Topics:
\begin{itemize}[leftmargin=2em]
    \item Circle vocabulary: radius, diameter, chord, secant, tangent
    \item Arcs and central angles
    \item Inscribed angles
    \item Angles inside and outside circles
    \item Segment relationships in circles
    \item Equations of circles: $(x-h)^2 + (y-k)^2 = r^2$
    \item Arc length: $s = r\theta$
    \item Sector area: $A = \tfrac{1}{2}r^2\theta$
\end{itemize}

\subsection*{Unit 10: Area, Surface Area, and Volume (5--6 weeks)}
Formulas:
\begin{itemize}[leftmargin=2em]
    \item Area of parallelogram: $A = bh$
    \item Area of triangle: $A = \tfrac{1}{2}bh$
    \item Area of trapezoid: $A = \tfrac{1}{2}(b_1 + b_2)h$
    \item Circle area: $A = \pi r^2$
    \item Prism volume: $V = Bh$
    \item Cylinder volume: $V = \pi r^2 h$
    \item Pyramid volume: $V = \tfrac{1}{3}Bh$
    \item Cone volume: $V = \tfrac{1}{3}\pi r^2 h$
    \item Sphere volume: $V = \tfrac{4}{3}\pi r^3$
    \item Sphere surface area: $SA = 4\pi r^2$
\end{itemize}

\section{Grading Breakdown}
\begin{center}
\begin{tabular}{L{0.55\linewidth}L{0.25\linewidth}}
\toprule
\textbf{Category} & \textbf{Percentage} \\
\midrule
Tests & 40\% \\
Quizzes & 25\% \\
Homework & 20\% \\
Projects/Constructions & 10\% \\
Class Participation & 5\% \\
\bottomrule
\end{tabular}
\end{center}

\chapter{Honors Geometry}

\noindent\textbf{HONORS COURSE:} Honors Geometry\\
\textbf{Grade Level:} 9th--10th Grade\\
\textbf{Duration:} Full Year (36 weeks)\\
\textbf{Credits:} 1.0 Credit (Weighted 4.5 scale)\\
\textbf{Prerequisites:} Honors Algebra 1 or teacher recommendation

\section{Course Description}
Honors Geometry extends the core geometry curriculum through deeper proof work, richer algebraic connections,
and more challenging applications. Students develop precision in formal reasoning while exploring
transformations, trigonometry, circles, and three-dimensional measurement at an accelerated pace.

\section{Distinctions from Geometry}
\begin{itemize}[leftmargin=2em]
    \item Greater emphasis on formal proof and justification
    \item More complex multi-step constructions and coordinate proofs
    \item Enrichment with loci, introductory vectors, and advanced circle relationships
    \item Faster-paced instruction with cumulative assessments and independent practice
\end{itemize}

\section{Units of Study}
\begin{itemize}[leftmargin=2em]
    \item Foundations of proof, logic, and constructions
    \item Congruence, transformations, and triangle relationships
    \item Similarity, right triangle trigonometry, and geometric modeling
    \item Circles, analytic geometry, and area/volume optimization
\end{itemize}

\section{Grading Breakdown}
\begin{center}
\begin{tabular}{L{0.55\linewidth}L{0.25\linewidth}}
\toprule
\textbf{Category} & \textbf{Percentage} \\
\midrule
Tests and Proof Assessments & 40\% \\
Quizzes & 20\% \\
Homework and Practice & 20\% \\
Projects and Constructions & 15\% \\
Class Participation & 5\% \\
\bottomrule
\end{tabular}
\end{center}

\chapter{Algebra 2}

\noindent\textbf{Course:} Algebra 2\\
\textbf{Grade Level:} 10th--11th Grade\\
\textbf{Duration:} Full Year (36 weeks)\\
\textbf{Credits:} 1.0 Credit\\
\textbf{Prerequisites:} Algebra 1 and Geometry

\section{Course Description}
Algebra 2 deepens students' understanding of functions, equations, and mathematical modeling. The course
emphasizes polynomial, rational, exponential, logarithmic, and trigonometric relationships while preparing
students for higher-level mathematics and STEM coursework.

\section{Course Goals}
\begin{itemize}[leftmargin=2em]
    \item Analyze and compare function families in multiple representations
    \item Solve nonlinear equations and systems using algebraic and graphical methods
    \item Model real-world situations with exponential, logarithmic, and polynomial functions
    \item Work fluently with complex numbers, sequences, and probability
\end{itemize}

\section{Units of Study}
\begin{itemize}[leftmargin=2em]
    \item Linear systems, matrices, and inequalities
    \item Quadratic, polynomial, and rational functions
    \item Radical, exponential, and logarithmic equations
    \item Sequences, series, probability, and data-informed decision making
\end{itemize}

\section{Grading Breakdown}
\begin{center}
\begin{tabular}{L{0.55\linewidth}L{0.25\linewidth}}
\toprule
\textbf{Category} & \textbf{Percentage} \\
\midrule
Tests & 35\% \\
Quizzes & 25\% \\
Homework & 20\% \\
Projects and Labs & 10\% \\
Final Exam and Participation & 10\% \\
\bottomrule
\end{tabular}
\end{center}

\chapter{Honors Algebra 2}

\noindent\textbf{HONORS COURSE:} Honors Algebra 2\\
\textbf{Grade Level:} 9th--11th Grade\\
\textbf{Duration:} Full Year (36 weeks)\\
\textbf{Credits:} 1.0 Credit (Weighted 4.5 scale)\\
\textbf{Prerequisites:} Geometry with teacher recommendation or Honors Geometry

\section{Course Description}
Honors Algebra 2 covers the full Algebra 2 curriculum with additional rigor, abstraction, and pace. Students
explore advanced function analysis, proof-based algebraic reasoning, and richer modeling tasks that prepare
them for Pre-Calculus, AP Precalculus, and AP Calculus pathways.

\section{Distinctions from Algebra 2}
\begin{itemize}[leftmargin=2em]
    \item Greater emphasis on symbolic manipulation and derivations
    \item Enrichment in conics, recursion, and advanced sequences and series
    \item Multi-representation analysis using algebraic, numeric, graphical, and verbal models
    \item More frequent cumulative assessments and non-routine problem solving
\end{itemize}

\section{Units of Study}
\begin{itemize}[leftmargin=2em]
    \item Advanced linear systems, matrices, and transformations
    \item Polynomial, rational, and radical function behavior
    \item Exponential and logarithmic modeling with applications
    \item Trigonometric foundations, conic sections, and probability/statistics
\end{itemize}

\section{Grading Breakdown}
\begin{center}
\begin{tabular}{L{0.55\linewidth}L{0.25\linewidth}}
\toprule
\textbf{Category} & \textbf{Percentage} \\
\midrule
Tests and Major Assessments & 40\% \\
Quizzes & 20\% \\
Homework and Practice & 20\% \\
Projects and Investigations & 15\% \\
Class Participation & 5\% \\
\bottomrule
\end{tabular}
\end{center}

\chapter{Pre-Calculus}

\noindent\textbf{Course:} Pre-Calculus\\
\textbf{Grade Level:} 11th--12th Grade\\
\textbf{Duration:} Full Year (36 weeks)\\
\textbf{Credits:} 1.0 Credit\\
\textbf{Prerequisites:} Algebra 2

\section{Course Description}
Pre-Calculus synthesizes algebra, geometry, and trigonometry into a unified study of advanced functions.
Students prepare for calculus by analyzing rates of change, functional composition, trigonometric models,
vectors, and polar representations.

\section{Course Goals}
\begin{itemize}[leftmargin=2em]
    \item Interpret and transform advanced function models
    \item Apply trigonometric identities and solve trigonometric equations
    \item Use vectors, parametric equations, and polar coordinates to model motion
    \item Strengthen readiness for limits, derivatives, and college mathematics
\end{itemize}

\section{Units of Study}
\begin{itemize}[leftmargin=2em]
    \item Function analysis, composition, and inverses
    \item Trigonometric functions, identities, and applications
    \item Complex numbers, vectors, matrices, and conics
    \item Parametric equations, polar graphs, and introductory limits
\end{itemize}

\section{Grading Breakdown}
\begin{center}
\begin{tabular}{L{0.55\linewidth}L{0.25\linewidth}}
\toprule
\textbf{Category} & \textbf{Percentage} \\
\midrule
Tests & 35\% \\
Quizzes & 20\% \\
Homework & 20\% \\
Projects and Modeling Tasks & 15\% \\
Final Exam and Participation & 10\% \\
\bottomrule
\end{tabular}
\end{center}

\chapter{Honors Pre-Calculus}

\noindent\textbf{HONORS COURSE:} Honors Pre-Calculus\\
\textbf{Grade Level:} 10th--12th Grade\\
\textbf{Duration:} Full Year (36 weeks)\\
\textbf{Credits:} 1.0 Credit (Weighted 4.5 scale)\\
\textbf{Prerequisites:} Honors Algebra 2 or teacher recommendation

\section{Course Description}
Honors Pre-Calculus accelerates students through advanced algebraic and trigonometric topics with a strong
focus on proof, abstraction, and calculus readiness. The course serves as a bridge to AP Precalculus,
AP Calculus AB, AP Calculus BC, and dual-credit calculus experiences.

\section{Distinctions from Pre-Calculus}
\begin{itemize}[leftmargin=2em]
    \item Deeper treatment of identities, inverse functions, and analytic trigonometry
    \item Advanced modeling with vectors, sequences, limits, and recursive reasoning
    \item Increased emphasis on justification, derivation, and technology-supported exploration
    \item More complex cumulative exams and independent problem sets
\end{itemize}

\section{Units of Study}
\begin{itemize}[leftmargin=2em]
    \item Advanced function transformations and proof of identities
    \item Trigonometric analysis, applications, and solving techniques
    \item Conic sections, vectors, matrices, and polar systems
    \item Sequences, series, introductory limits, and calculus-aligned modeling
\end{itemize}

\section{Grading Breakdown}
\begin{center}
\begin{tabular}{L{0.55\linewidth}L{0.25\linewidth}}
\toprule
\textbf{Category} & \textbf{Percentage} \\
\midrule
Tests and Major Assessments & 40\% \\
Quizzes & 20\% \\
Homework and Practice & 20\% \\
Projects and Investigations & 15\% \\
Class Participation & 5\% \\
\bottomrule
\end{tabular}
\end{center}

\chapter{AP Precalculus}

\noindent\textbf{AP COURSE:} AP Precalculus\\
\textbf{Grade Level:} 10th--12th Grade\\
\textbf{Duration:} Full Year (36 weeks)\\
\textbf{Credits:} 1.0 Credit (Weighted 5.0 scale)\\
\textbf{Prerequisites:} Honors Algebra 2 or strong performance in Algebra 2 with recommendation

\section{Course Description}
AP Precalculus develops students' mastery of mathematical functions, modeling, and multiple representations
through the College Board framework. The course emphasizes conceptual understanding, procedural fluency,
and real-world applications in preparation for AP assessment and future STEM coursework.

\section{AP Course Framework}
\begin{itemize}[leftmargin=2em]
    \item Polynomial and rational functions
    \item Exponential and logarithmic functions
    \item Trigonometric and polar functions
    \item Functions involving parameters, vectors, and dynamic modeling
\end{itemize}

\section{AP Skills Emphasized}
\begin{itemize}[leftmargin=2em]
    \item Translating among symbolic, graphical, numerical, and verbal representations
    \item Justifying procedures and interpreting results in context
    \item Modeling authentic scenarios using technology strategically
    \item Preparing for the AP Precalculus Exam through timed practice and performance tasks
\end{itemize}

\section{Grading Breakdown}
\begin{center}
\begin{tabular}{L{0.55\linewidth}L{0.25\linewidth}}
\toprule
\textbf{Category} & \textbf{Percentage} \\
\midrule
Unit Tests and AP-Style Assessments & 35\% \\
Quizzes & 20\% \\
Homework and Practice & 15\% \\
Projects and Performance Tasks & 20\% \\
AP Exam Preparation & 10\% \\
\bottomrule
\end{tabular}
\end{center}

\chapter{Calculus I (Dual Credit)}

\noindent\textbf{DUAL CREDIT COURSE:} Calculus I (Dual Credit)\\
\textbf{Grade Level:} 11th--12th Grade\\
\textbf{Duration:} One Semester or Full Year Equivalent\\
\textbf{Credits:} High School Credit + College Credit\\
\textbf{Prerequisites:} Pre-Calculus and qualifying placement scores

\section{Course Description}
Calculus I (Dual Credit) introduces differential and introductory integral calculus through a college-aligned
syllabus. Students earn both high school and transferable college credit while developing fluency with limits,
derivatives, optimization, related rates, and the Fundamental Theorem of Calculus.

\section{College-Readiness Expectations}
\begin{itemize}[leftmargin=2em]
    \item Independent note-taking, study planning, and assignment completion
    \item Strong algebra and trigonometry fluency without extensive remediation
    \item Participation in college-style assessments and cumulative exams
    \item Clear communication of mathematical reasoning in written work
\end{itemize}

\section{Units of Study}
\begin{itemize}[leftmargin=2em]
    \item Limits, continuity, and introductory proof ideas
    \item Derivatives and differentiation techniques
    \item Applications of derivatives including optimization and related rates
    \item Antiderivatives, definite integrals, and the Fundamental Theorem of Calculus
\end{itemize}

\section{Grading Breakdown}
\begin{center}
\begin{tabular}{L{0.55\linewidth}L{0.25\linewidth}}
\toprule
\textbf{Category} & \textbf{Percentage} \\
\midrule
Exams & 45\% \\
Quizzes & 20\% \\
Homework and Online Practice & 15\% \\
Projects or Labs & 10\% \\
Final Exam & 10\% \\
\bottomrule
\end{tabular}
\end{center}

\chapter{AP Calculus AB}

\noindent\textbf{AP COURSE:} AP Calculus AB\\
\textbf{Grade Level:} 11th--12th Grade\\
\textbf{Duration:} Full Year (36 weeks)\\
\textbf{Credits:} 1.0 Credit (Weighted 5.0 scale)\\
\textbf{Prerequisites:} Pre-Calculus with teacher recommendation

\section{Course Description}
AP Calculus AB is a first-semester college calculus course built around limits, derivatives, and integrals.
Students analyze functions, model change, and solve applied problems while preparing for the AP Calculus AB
Exam.

\section{AP Units}
\begin{itemize}[leftmargin=2em]
    \item Limits and continuity
    \item Differentiation: definition, rules, and contextual interpretation
    \item Applications of derivatives
    \item Integration and accumulation of change
    \item Differential equations and slope fields
\end{itemize}

\section{AP Skills Emphasized}
\begin{itemize}[leftmargin=2em]
    \item Connecting multiple representations of functions and rates of change
    \item Justifying solutions without unsupported calculator dependence
    \item Writing clear, College Board-aligned free-response explanations
    \item Building stamina for timed multiple-choice and free-response sections
\end{itemize}

\section{Grading Breakdown}
\begin{center}
\begin{tabular}{L{0.55\linewidth}L{0.25\linewidth}}
\toprule
\textbf{Category} & \textbf{Percentage} \\
\midrule
Unit Tests and Free-Response Sets & 40\% \\
Quizzes & 20\% \\
Homework and Practice & 15\% \\
Labs and Investigations & 10\% \\
AP Exam Preparation & 15\% \\
\bottomrule
\end{tabular}
\end{center}

\chapter{AP Calculus BC}

\noindent\textbf{AP COURSE:} AP Calculus BC\\
\textbf{Grade Level:} 11th--12th Grade\\
\textbf{Duration:} Full Year (36 weeks)\\
\textbf{Credits:} 1.0 Credit (Weighted 5.0 scale)\\
\textbf{Prerequisites:} Honors Pre-Calculus or AP Precalculus with recommendation

\section{Course Description}
AP Calculus BC includes the entire AP Calculus AB curriculum and extends it through sequences and series,
parametric, polar, and vector-valued functions, and advanced integration techniques. The course is designed
for students pursuing highly rigorous STEM pathways.

\section{BC-Specific Extensions}
\begin{itemize}[leftmargin=2em]
    \item Parametric, polar, and vector-valued motion
    \item Integration by parts, partial fractions, and improper integrals
    \item Logistic growth and richer differential equation applications
    \item Convergence tests and Taylor/Maclaurin polynomial approximations
\end{itemize}

\section{AP Skills Emphasized}
\begin{itemize}[leftmargin=2em]
    \item Managing complex multi-step reasoning under timed conditions
    \item Connecting series representations to function behavior and approximation
    \item Interpreting results across graphical, numeric, algebraic, and contextual forms
    \item Preparing for both AB subscore and BC exam expectations
\end{itemize}

\section{Grading Breakdown}
\begin{center}
\begin{tabular}{L{0.55\linewidth}L{0.25\linewidth}}
\toprule
\textbf{Category} & \textbf{Percentage} \\
\midrule
Unit Tests and Free-Response Sets & 40\% \\
Quizzes & 20\% \\
Homework and Practice & 15\% \\
Labs and Investigations & 10\% \\
AP Exam Preparation & 15\% \\
\bottomrule
\end{tabular}
\end{center}

\chapter{Statistics}

\noindent\textbf{Course:} Statistics\\
\textbf{Grade Level:} 11th--12th Grade\\
\textbf{Duration:} Full Year (36 weeks)\\
\textbf{Credits:} 1.0 Credit\\
\textbf{Prerequisites:} Algebra 2

\section{Course Description}
Statistics introduces students to collecting, analyzing, and interpreting data in real contexts. The course
focuses on study design, probability, distributions, inference, and communication of conclusions using
graphs, summaries, and statistical reasoning.

\section{Course Goals}
\begin{itemize}[leftmargin=2em]
    \item Summarize univariate and bivariate data accurately
    \item Use probability models to make predictions and decisions
    \item Understand sampling methods, bias, and experimental design
    \item Draw conclusions using confidence intervals and significance tests
\end{itemize}

\section{Units of Study}
\begin{itemize}[leftmargin=2em]
    \item Describing distributions and comparing data sets
    \item Correlation, regression, and modeling relationships
    \item Probability, random variables, and sampling distributions
    \item Statistical inference for proportions and means
\end{itemize}

\section{Grading Breakdown}
\begin{center}
\begin{tabular}{L{0.55\linewidth}L{0.25\linewidth}}
\toprule
\textbf{Category} & \textbf{Percentage} \\
\midrule
Tests & 35\% \\
Quizzes & 20\% \\
Homework & 15\% \\
Projects and Data Investigations & 20\% \\
Class Participation and Labs & 10\% \\
\bottomrule
\end{tabular}
\end{center}

\chapter{AP Statistics}

\noindent\textbf{AP COURSE:} AP Statistics\\
\textbf{Grade Level:} 11th--12th Grade\\
\textbf{Duration:} Full Year (36 weeks)\\
\textbf{Credits:} 1.0 Credit (Weighted 5.0 scale)\\
\textbf{Prerequisites:} Algebra 2 with strong analytical skills

\section{Course Description}
AP Statistics is an introductory college-level statistics course following the College Board framework.
Students study exploratory data analysis, probability, sampling, inference, and experimental design while
communicating conclusions with precision and context awareness.

\section{AP Course Themes}
\begin{itemize}[leftmargin=2em]
    \item Exploring one-variable and two-variable data
    \item Collecting data through sampling and experimentation
    \item Anticipating patterns through probability and simulation
    \item Statistical inference using confidence intervals and tests
\end{itemize}

\section{AP Skills Emphasized}
\begin{itemize}[leftmargin=2em]
    \item Writing justified conclusions that reference statistical evidence
    \item Interpreting significance, confidence, bias, and variability in context
    \item Using graphing technology and simulation strategically
    \item Preparing for AP multiple-choice and free-response tasks
\end{itemize}

\section{Grading Breakdown}
\begin{center}
\begin{tabular}{L{0.55\linewidth}L{0.25\linewidth}}
\toprule
\textbf{Category} & \textbf{Percentage} \\
\midrule
Unit Tests and AP-Style Assessments & 35\% \\
Quizzes & 20\% \\
Homework and Practice & 15\% \\
Projects and Investigations & 20\% \\
AP Exam Preparation & 10\% \\
\bottomrule
\end{tabular}
\end{center}

% =========================
% Part III: Sciences
% =========================
\part{Sciences}

\chapter{Biology}

\noindent\textbf{Course:} Biology\\
\textbf{Grade Level:} 9th--10th Grade\\
\textbf{Duration:} Full Year (36 weeks)\\
\textbf{Credits:} 1.0 Credit\\
\textbf{Prerequisites:} Physical Science or equivalent (recommended)

\section{Course Description}
Biology is a laboratory-based life science course that introduces students to fundamental concepts
in cell biology, genetics, evolution, ecology, and human body systems. Emphasis is placed on
scientific inquiry, experimental design, and data analysis using real-world biological phenomena.

\section{Course Goals}
Upon completion of this course, students will be able to:
\begin{itemize}[leftmargin=2em]
    \item Apply the scientific method and laboratory skills to investigate biological questions
    \item Explain the structure and function of cells and cellular processes
    \item Describe patterns of inheritance and the molecular basis of genetics
    \item Analyze evidence for biological evolution and common ancestry
    \item Describe the flow of energy and matter in ecosystems
    \item Explain major human body systems and homeostasis
    \item Interpret and communicate scientific data using appropriate terminology
\end{itemize}

\section{Required Materials}
\begin{itemize}[leftmargin=2em]
    \item Laboratory notebook or bound composition book
    \item Scientific or graphing calculator
    \item Safety goggles (school-issued or personal, as required)
    \item Three-ring binder with dividers and loose-leaf paper
    \item Colored pencils for diagrams
\end{itemize}

\section{Laboratory Safety}
Students must:
\begin{itemize}[leftmargin=2em]
    \item Follow all written and verbal safety instructions
    \item Wear goggles, gloves, and aprons when required
    \item Handle glassware, chemicals, and biological specimens responsibly
    \item Never eat or drink in the laboratory
    \item Report all spills, accidents, or injuries immediately
    \item Properly dispose of chemical and biological wastes
\end{itemize}
Failure to follow safety procedures may result in removal from the lab activity and a loss of credit.

\section{Units of Study}
\subsection*{Unit 1: Introduction to Biology and Scientific Skills}
\begin{itemize}[leftmargin=2em]
    \item Nature of science and scientific method
    \item Experimental design and controlled experiments
    \item Measurement, data collection, and graphing
    \item Use of microscopes and basic lab equipment
\end{itemize}

\subsection*{Unit 2: Biochemistry}
\begin{itemize}[leftmargin=2em]
    \item Structure of atoms and molecules
    \item Properties of water and pH
    \item Organic molecules: carbohydrates, lipids, proteins, nucleic acids
    \item Enzymes and reaction rates
\end{itemize}

\subsection*{Unit 3: Cell Structure and Function}
\begin{itemize}[leftmargin=2em]
    \item Prokaryotic vs.\ eukaryotic cells
    \item Organelles and their functions
    \item Cell membrane structure and transport
    \item Osmosis, diffusion, and active transport
    \item Cell communication and homeostasis
\end{itemize}

\subsection*{Unit 4: Energy in Cells}
\begin{itemize}[leftmargin=2em]
    \item Photosynthesis: light-dependent and light-independent reactions
    \item Cellular respiration: glycolysis, Krebs cycle, electron transport chain
    \item ATP and energy transfer
\end{itemize}

\subsection*{Unit 5: Cell Division and Genetics}
\begin{itemize}[leftmargin=2em]
    \item Cell cycle, mitosis, and meiosis
    \item Mendelian genetics and Punnett squares
    \item Complex inheritance patterns
    \item DNA structure and replication
    \item Protein synthesis (transcription and translation)
    \item Mutations and genetic technology (overview)
\end{itemize}

\subsection*{Unit 6: Evolution}
\begin{itemize}[leftmargin=2em]
    \item Evidence for evolution
    \item Natural selection and adaptation
    \item Speciation and phylogenetic trees
    \item Population genetics (basic)
\end{itemize}

\subsection*{Unit 7: Ecology}
\begin{itemize}[leftmargin=2em]
    \item Ecosystem organization (organism to biosphere)
    \item Food chains, food webs, and energy pyramids
    \item Biogeochemical cycles
    \item Population dynamics and carrying capacity
    \item Human impact on the environment
\end{itemize}

\subsection*{Unit 8: Human Body Systems (Survey)}
\begin{itemize}[leftmargin=2em]
    \item Major organ systems and homeostasis
    \item Interactions among systems
\end{itemize}

\section{Grading Breakdown}
\begin{center}
\begin{tabular}{L{0.55\linewidth}L{0.25\linewidth}}
\toprule
\textbf{Category} & \textbf{Percentage} \\
\midrule
Tests and Major Assessments & 35\% \\
Laboratory Reports and Practical Exams & 25\% \\
Quizzes & 15\% \\
Classwork and Homework & 15\% \\
Participation and Lab Conduct & 10\% \\
\bottomrule
\end{tabular}
\end{center}

\chapter{Honors Biology}

\noindent\textbf{HONORS COURSE:} Honors Biology\\
\textbf{Grade Level:} 9th--10th Grade\\
\textbf{Duration:} Full Year (36 weeks)\\
\textbf{Credits:} 1.0 Credit (Weighted 4.5 scale)\\
\textbf{Prerequisites:} Strong performance in previous science course and teacher recommendation

\section{Course Description}
Honors Biology covers all topics of Biology at an accelerated pace with greater depth, more rigorous
laboratory work, and increased emphasis on scientific reading and writing. The course is designed
to prepare students for AP Biology and other advanced science coursework.

\section{Distinctions from Regular Biology}
\begin{itemize}[leftmargin=2em]
    \item Faster pacing and more complex content
    \item Additional and more sophisticated laboratory investigations
    \item Higher expectations for independent reading of scientific texts
    \item More quantitative data analysis and graphing
    \item Extended inquiry projects and research-style reports
\end{itemize}

\section{Additional Topics}
\begin{itemize}[leftmargin=2em]
    \item Molecular genetics in greater depth
    \item Population genetics and Hardy--Weinberg equilibrium (introductory)
    \item Biostatistics and error analysis in labs
    \item Current issues in biotechnology and bioethics
\end{itemize}

\chapter{AP Biology}

\noindent\textbf{AP COURSE:} AP Biology\\
\textbf{Grade Level:} 11th--12th Grade\\
\textbf{Duration:} Full Year (36 weeks)\\
\textbf{Credits:} 1.0 Credit (Weighted 5.0 scale)\\
\textbf{Prerequisites:} Biology and Chemistry (Honors recommended) and teacher recommendation

\section{Course Description}
AP Biology is a college-level biology course that follows the College Board AP Biology framework.
Students engage in in-depth study of biological concepts, emphasizing inquiry-based laboratory
investigations, data analysis, and application of scientific practices. The course prepares students
for the AP Biology exam in May.

\section{AP Exam Components}
\subsection*{Multiple Choice and Grid-In (50\% of score)}
\begin{itemize}[leftmargin=2em]
    \item 60 multiple-choice questions
    \item 6 grid-in (quantitative) questions
\end{itemize}

\subsection*{Free Response (50\% of score)}
\begin{itemize}[leftmargin=2em]
    \item 2 long free-response questions
    \item 4 short free-response questions
    \item Emphasis on experimental design, data analysis, and conceptual understanding
\end{itemize}

\section{Big Ideas (Enduring Understandings)}
\begin{itemize}[leftmargin=2em]
    \item \textbf{Big Idea 1:} Evolution drives the diversity and unity of life.
    \item \textbf{Big Idea 2:} Biological systems utilize energy and molecular building blocks.
    \item \textbf{Big Idea 3:} Living systems store, retrieve, transmit, and respond to information.
    \item \textbf{Big Idea 4:} Biological systems interact and exhibit complex properties.
\end{itemize}

\section{Science Practices}
Students develop proficiency in:
\begin{itemize}[leftmargin=2em]
    \item Concept explanation
    \item Visual representation
    \item Questions and methods
    \item Data collection and analysis
    \item Statistical tests and data analysis
    \item Argumentation and justification
\end{itemize}

\section{Units of Study (College Board Framework)}
\subsection*{Unit 1: Chemistry of Life}
\subsection*{Unit 2: Cell Structure and Function}
\subsection*{Unit 3: Cellular Energetics}
\subsection*{Unit 4: Cell Communication and Cell Cycle}
\subsection*{Unit 5: Heredity}
\subsection*{Unit 6: Gene Expression and Regulation}
\subsection*{Unit 7: Natural Selection}
\subsection*{Unit 8: Ecology}

Each unit integrates required AP Biology labs, data analysis, and practice with AP-style questions.

\section{Required Laboratory Investigations}
Students complete a series of inquiry-based labs such as:
\begin{itemize}[leftmargin=2em]
    \item Diffusion and osmosis
    \item Enzyme activity
    \item Photosynthesis and cellular respiration
    \item Cell division and genetics
    \item Population genetics and natural selection
    \item Energy dynamics in ecosystems
\end{itemize}

\section{Grading Breakdown}
\begin{center}
\begin{tabular}{L{0.55\linewidth}L{0.25\linewidth}}
\toprule
\textbf{Category} & \textbf{Percentage} \\
\midrule
Unit Tests and AP-Style Exams & 40\% \\
Laboratory Reports and Practicals & 30\% \\
Quizzes and Problem Sets & 15\% \\
Homework and Reading Checks & 10\% \\
Participation and Lab Practices & 5\% \\
\bottomrule
\end{tabular}
\end{center}

\chapter{Chemistry}

\noindent\textbf{Course:} Chemistry\\
\textbf{Grade Level:} 10th--11th Grade\\
\textbf{Duration:} Full Year (36 weeks)\\
\textbf{Credits:} 1.0 Credit\\
\textbf{Prerequisites:} Algebra 1 and Biology (concurrent enrollment may be allowed)

\section{Course Description}
Chemistry is a quantitative laboratory science focused on the composition, structure, properties,
and reactions of matter. Students develop strong problem-solving skills while investigating atomic
structure, bonding, chemical reactions, stoichiometry, states of matter, and basic thermochemistry.

\section{Course Goals}
Students will:
\begin{itemize}[leftmargin=2em]
    \item Use scientific measurement and dimensional analysis
    \item Explain atomic structure and periodic trends
    \item Describe different types of chemical bonding
    \item Predict and balance chemical reactions
    \item Perform stoichiometric calculations
    \item Analyze gas behavior and solution properties
    \item Safely conduct laboratory experiments and interpret data
\end{itemize}

\section{Required Materials}
\begin{itemize}[leftmargin=2em]
    \item Scientific or graphing calculator
    \item Laboratory notebook
    \item Safety goggles (as required)
    \item Three-ring binder and loose-leaf paper
\end{itemize}

\section{Units of Study}
\subsection*{Unit 1: Scientific Measurement and Safety}
\subsection*{Unit 2: Matter and Energy}
\subsection*{Unit 3: Atomic Structure and the Periodic Table}
\subsection*{Unit 4: Chemical Bonding and Nomenclature}
\subsection*{Unit 5: Chemical Reactions and Stoichiometry}
\subsection*{Unit 6: States of Matter and Gas Laws}
\subsection*{Unit 7: Solutions, Acids, and Bases}
\subsection*{Unit 8: Thermochemistry (Introductory)}

\section{Grading Breakdown}
\begin{center}
\begin{tabular}{L{0.55\linewidth}L{0.25\linewidth}}
\toprule
\textbf{Category} & \textbf{Percentage} \\
\midrule
Tests and Major Assessments & 35\% \\
Laboratory Reports and Practicals & 25\% \\
Quizzes & 20\% \\
Homework and Classwork & 15\% \\
Participation and Safety & 5\% \\
\bottomrule
\end{tabular}
\end{center}

\chapter{Honors Chemistry}

\noindent\textbf{HONORS COURSE:} Honors Chemistry\\
\textbf{Grade Level:} 10th--11th Grade\\
\textbf{Duration:} Full Year (36 weeks)\\
\textbf{Credits:} 1.0 Credit (Weighted 4.5 scale)\\
\textbf{Prerequisites:} Honors Geometry or Algebra 2 (concurrent) and teacher recommendation

\section{Course Description}
Honors Chemistry is an accelerated and more mathematically rigorous version of Chemistry.
Students study the same core topics with additional depth, extended quantitative problem-solving,
and more complex laboratory investigations, preparing them for AP Chemistry and other advanced
science courses.

\section{Distinctions from Regular Chemistry}
\begin{itemize}[leftmargin=2em]
    \item Greater emphasis on algebraic and proportional reasoning
    \item More challenging stoichiometry and gas law problems
    \item Introduction to equilibrium and thermodynamics concepts
    \item More frequent and in-depth laboratory investigations
    \item AP-style free-response and multiple-choice practice questions
\end{itemize}

\section{Additional Topics}
\begin{itemize}[leftmargin=2em]
    \item Basic chemical equilibrium
    \item Introductory kinetics
    \item Acid--base titrations and pH calculations
    \item Thermochemical equations and calorimetry
\end{itemize}

\chapter{AP Chemistry}

\noindent\textbf{AP COURSE:} AP Chemistry\\
\textbf{Grade Level:} 11th--12th Grade\\
\textbf{Duration:} Full Year (36 weeks)\\
\textbf{Credits:} 1.0 Credit (Weighted 5.0 scale)\\
\textbf{Prerequisites:} Chemistry (Honors recommended) and concurrent enrollment in Pre-Calculus or higher

\section{Course Description}
AP Chemistry is a college-level course following the College Board AP Chemistry curriculum
framework. The course emphasizes inquiry-based labs, conceptual understanding, and quantitative
problem-solving across all major areas of general chemistry.

\section{AP Exam Components}
\begin{itemize}[leftmargin=2em]
    \item Multiple-choice section (MCQ)
    \item Free-response section including long and short questions
    \item Emphasis on explanation, representation, and calculations
\end{itemize}

\section{Units of Study (College Board Framework)}
\subsection*{Unit 1: Atomic Structure and Properties}
\subsection*{Unit 2: Molecular and Ionic Compound Structure and Properties}
\subsection*{Unit 3: Intermolecular Forces and Properties}
\subsection*{Unit 4: Chemical Reactions}
\subsection*{Unit 5: Kinetics}
\subsection*{Unit 6: Thermodynamics}
\subsection*{Unit 7: Equilibrium}
\subsection*{Unit 8: Acids and Bases}
\subsection*{Unit 9: Applications of Thermodynamics}

\section{Required Laboratory Investigations}
Students perform a set of college-level labs addressing:
\begin{itemize}[leftmargin=2em]
    \item Stoichiometry and limiting reactants
    \item Titrations and pH measurements
    \item Gas laws and molar mass determination
    \item Enthalpy changes and calorimetry
    \item Equilibrium constants and Le Châtelier's principle
\end{itemize}

\section{Grading Breakdown}
\begin{center}
\begin{tabular}{L{0.55\linewidth}L{0.25\linewidth}}
\toprule
\textbf{Category} & \textbf{Percentage} \\
\midrule
Unit Exams and AP-Style Tests & 45\% \\
Laboratory Reports and Practicals & 30\% \\
Quizzes and Problem Sets & 15\% \\
Homework and Participation & 10\% \\
\bottomrule
\end{tabular}
\end{center}

\chapter{Physics}

\noindent\textbf{Course:} Physics\\
\textbf{Grade Level:} 11th--12th Grade\\
\textbf{Duration:} Full Year (36 weeks)\\
\textbf{Credits:} 1.0 Credit\\
\textbf{Prerequisites:} Algebra 2 (concurrent Pre-Calculus recommended)

\section{Course Description}
Physics is a laboratory-based course focusing on the study of motion, forces, energy, waves, and
electricity. Students use mathematical models and experimental data to describe and predict
physical phenomena.

\section{Course Goals}
\begin{itemize}[leftmargin=2em]
    \item Apply Newtonian mechanics to motion in one and two dimensions
    \item Analyze energy, work, and power in systems
    \item Investigate mechanical waves and sound
    \item Explore basic electricity and magnetism concepts
    \item Collect, analyze, and interpret experimental data
\end{itemize}

\section{Units of Study}
\subsection*{Unit 1: Kinematics}
\subsection*{Unit 2: Forces and Newton's Laws}
\subsection*{Unit 3: Work, Energy, and Power}
\subsection*{Unit 4: Momentum and Collisions}
\subsection*{Unit 5: Circular Motion and Gravitation}
\subsection*{Unit 6: Waves and Sound}
\subsection*{Unit 7: Electricity (Basic Circuits)}

\section{Grading Breakdown}
\begin{center}
\begin{tabular}{L{0.55\linewidth}L{0.25\linewidth}}
\toprule
\textbf{Category} & \textbf{Percentage} \\
\midrule
Tests & 40\% \\
Laboratory Reports and Practicals & 25\% \\
Quizzes & 15\% \\
Homework and Classwork & 15\% \\
Participation and Safety & 5\% \\
\bottomrule
\end{tabular}
\end{center}

\chapter{Honors Physics}

\noindent\textbf{HONORS COURSE:} Honors Physics\\
\textbf{Grade Level:} 11th--12th Grade\\
\textbf{Duration:} Full Year (36 weeks)\\
\textbf{Credits:} 1.0 Credit (Weighted 4.5 scale)\\
\textbf{Prerequisites:} Honors Algebra 2 or Pre-Calculus and teacher recommendation

\section{Course Description}
Honors Physics covers the same major topics as Physics with more mathematical rigor,
greater depth, and additional problem-solving. It serves as a strong foundation for
AP Physics courses.

\section{Distinctions from Regular Physics}
\begin{itemize}[leftmargin=2em]
    \item Use of more advanced algebra and trigonometry
    \item Multi-step, open-ended problems
    \item More demanding laboratory investigations
    \item Introduction to rotational dynamics and simple harmonic motion (as time permits)
\end{itemize}

\section{Additional Topics}
\begin{itemize}[leftmargin=2em]
    \item Vector analysis in two dimensions
    \item Rotational motion and torque
    \item Simple harmonic motion and resonance
\end{itemize}

\chapter{AP Physics 1}

\noindent\textbf{AP COURSE:} AP Physics 1: Algebra-Based\\
\textbf{Grade Level:} 11th--12th Grade\\
\textbf{Duration:} Full Year (36 weeks)\\
\textbf{Credits:} 1.0 Credit (Weighted 5.0 scale)\\
\textbf{Prerequisites:} Honors Physics (recommended) and concurrent enrollment in Algebra 2 or higher

\section{Course Description}
AP Physics 1 is an algebra-based, college-level physics course covering Newtonian mechanics,
mechanical waves, and introductory circuits. The course aligns with the College Board framework
and emphasizes conceptual understanding, quantitative reasoning, and inquiry-based labs.

\section{AP Exam Components}
\begin{itemize}[leftmargin=2em]
    \item Multiple-choice questions (single-select and multi-select)
    \item Free-response questions including experimental design and qualitative/quantitative reasoning
\end{itemize}

\section{Units of Study}
\subsection*{Unit 1: Kinematics}
\subsection*{Unit 2: Dynamics (Forces)}
\subsection*{Unit 3: Circular Motion and Gravitation}
\subsection*{Unit 4: Energy}
\subsection*{Unit 5: Momentum}
\subsection*{Unit 6: Simple Harmonic Motion}
\subsection*{Unit 7: Torque and Rotational Motion}
\subsection*{Unit 8: Electric Charge and DC Circuits (Introductory)}

\section{Science Practices}
Students develop:
\begin{itemize}[leftmargin=2em]
    \item Experimental design and data collection skills
    \item Graphical and mathematical modeling abilities
    \item Written explanations of physical reasoning
\end{itemize}

\section{Grading Breakdown}
\begin{center}
\begin{tabular}{L{0.55\linewidth}L{0.25\linewidth}}
\toprule
\textbf{Category} & \textbf{Percentage} \\
\midrule
AP-Style Exams and Unit Tests & 45\% \\
Laboratory Investigations and Reports & 30\% \\
Quizzes and Problem Sets & 15\% \\
Homework and Participation & 10\% \\
\bottomrule
\end{tabular}
\end{center}

\chapter{AP Physics 2}

\noindent\textbf{AP COURSE:} AP Physics 2: Algebra-Based\\
\textbf{Grade Level:} 11th--12th Grade\\
\textbf{Duration:} Full Year (36 weeks)\\
\textbf{Credits:} 1.0 Credit (Weighted 5.0 scale)\\
\textbf{Prerequisites:} AP Physics 1 or equivalent and teacher recommendation

\section{Course Description}
AP Physics 2 is an algebra-based, college-level physics course focusing on fluids, thermodynamics,
electricity and magnetism, optics, and modern physics, following the College Board framework.

\section{Units of Study}
\subsection*{Unit 1: Fluids}
\subsection*{Unit 2: Thermodynamics}
\subsection*{Unit 3: Electric Force, Fields, and Potential}
\subsection*{Unit 4: Electric Circuits}
\subsection*{Unit 5: Magnetism and Electromagnetic Induction}
\subsection*{Unit 6: Geometric and Physical Optics}
\subsection*{Unit 7: Quantum, Atomic, and Nuclear Physics}

\chapter{AP Physics C: Mechanics}

\noindent\textbf{AP COURSE:} AP Physics C: Mechanics\\
\textbf{Grade Level:} 12th Grade\\
\textbf{Duration:} Full Year or Semester (local decision)\\
\textbf{Credits:} 1.0 Credit or 0.5 Credit (Weighted 5.0 scale)\\
\textbf{Prerequisites:} Concurrent enrollment in Calculus AB/BC and prior physics course

\section{Course Description}
AP Physics C: Mechanics is a calculus-based, college-level course emphasizing rigorous
treatment of Newtonian mechanics using differential and integral calculus.

\section{AP Exam Components}
\begin{itemize}[leftmargin=2em]
    \item Multiple-choice section
    \item Free-response section with quantitative and qualitative problems
\end{itemize}

\section{Units of Study}
\subsection*{Unit 1: Kinematics (Calculus-Based)}
\subsection*{Unit 2: Newton's Laws of Motion}
\subsection*{Unit 3: Work, Energy, and Power}
\subsection*{Unit 4: Systems of Particles and Linear Momentum}
\subsection*{Unit 5: Circular Motion and Rotation}
\subsection*{Unit 6: Oscillations and Gravitation}

\section{Calculus Applications}
\begin{itemize}[leftmargin=2em]
    \item Using derivatives to find velocity and acceleration
    \item Using integrals to find displacement and work
    \item Differential equations in simple harmonic motion
\end{itemize}

\section{Grading Breakdown}
\begin{center}
\begin{tabular}{L{0.55\linewidth}L{0.25\linewidth}}
\toprule
\textbf{Category} & \textbf{Percentage} \\
\midrule
Exams and AP-Style Tests & 50\% \\
Problem Sets & 25\% \\
Laboratory Investigations & 15\% \\
Homework and Participation & 10\% \\
\bottomrule
\end{tabular}
\end{center}

\chapter{AP Physics C: Electricity and Magnetism}

\noindent\textbf{AP COURSE:} AP Physics C: Electricity and Magnetism\\
\textbf{Grade Level:} 12th Grade\\
\textbf{Duration:} Semester or Full Year (paired with Mechanics)\\
\textbf{Credits:} 0.5--1.0 Credit (Weighted 5.0 scale)\\
\textbf{Prerequisites:} AP Physics C: Mechanics (recommended) and concurrent Calculus

\section{Course Description}
AP Physics C: Electricity and Magnetism is a calculus-based course exploring electrostatics,
conductors, capacitors, dielectrics, electric circuits, magnetism, and electromagnetic induction.

\section{Units of Study}
\subsection*{Unit 1: Electrostatics}
\subsection*{Unit 2: Conductors, Capacitors, and Dielectrics}
\subsection*{Unit 3: Electric Circuits}
\subsection*{Unit 4: Magnetic Fields}
\subsection*{Unit 5: Electromagnetic Induction}

% =========================
% Part IV: Social Studies and History
% =========================
\part{Social Studies and History}

\chapter{World History}

\noindent\textbf{Course:} World History\\
\textbf{Grade Level:} 10th Grade\\
\textbf{Duration:} Full Year (36 weeks)\\
\textbf{Credits:} 1.0 Credit\\
\textbf{Prerequisites:} None

\section{Course Description}
World History surveys major civilizations, events, and themes from ancient times to the present.
The course emphasizes chronological reasoning, comparative analysis, and historical thinking skills.

\section{Course Goals}
\begin{itemize}[leftmargin=2em]
    \item Trace major developments in world civilizations over time
    \item Analyze primary and secondary sources
    \item Compare political, economic, and social systems
    \item Develop historical arguments supported by evidence
\end{itemize}

\section{Units of Study}
\subsection*{Unit 1: Foundations of Civilization}
\subsection*{Unit 2: Classical Civilizations}
\subsection*{Unit 3: Postclassical Era and Global Interactions}
\subsection*{Unit 4: Early Modern World}
\subsection*{Unit 5: Revolutions and Industrialization}
\subsection*{Unit 6: World Wars and Global Conflict}
\subsection*{Unit 7: Contemporary World}

\section{Historical Thinking Skills}
Students practice:
\begin{itemize}[leftmargin=2em]
    \item Sourcing and contextualization
    \item Causation and continuity/change over time
    \item Comparison
    \item Argumentation and use of evidence
\end{itemize}

\section{Grading Breakdown}
\begin{center}
\begin{tabular}{L{0.55\linewidth}L{0.25\linewidth}}
\toprule
\textbf{Category} & \textbf{Percentage} \\
\midrule
Tests and DBQ/Essay Assessments & 40\% \\
Quizzes & 20\% \\
Projects and Presentations & 20\% \\
Homework and Classwork & 15\% \\
Participation & 5\% \\
\bottomrule
\end{tabular}
\end{center}

\chapter{AP World History: Modern}

\noindent\textbf{AP COURSE:} AP World History: Modern\\
\textbf{Grade Level:} 10th--11th Grade\\
\textbf{Duration:} Full Year (36 weeks)\\
\textbf{Credits:} 1.0 Credit (Weighted 5.0 scale)\\
\textbf{Prerequisites:} Strong reading and writing skills; teacher recommendation

\section{Course Description}
AP World History: Modern is a college-level course focusing on world history from c.\ 1200 CE
to the present. Students develop advanced historical thinking skills and prepare for the AP exam.

\section{AP Exam Components}
\begin{itemize}[leftmargin=2em]
    \item Multiple-choice questions on stimulus-based sets
    \item Short-answer questions (SAQs)
    \item Document-Based Question (DBQ)
    \item Long Essay Question (LEQ)
\end{itemize}

\section{Units of Study}
\subsection*{Unit 1: The Global Tapestry (c.\ 1200--c.\ 1450)}
\subsection*{Unit 2: Networks of Exchange}
\subsection*{Unit 3: Land-Based Empires}
\subsection*{Unit 4: Transoceanic Interconnections}
\subsection*{Unit 5: Revolutions}
\subsection*{Unit 6: Consequences of Industrialization}
\subsection*{Unit 7: Global Conflict}
\subsection*{Unit 8: Cold War and Decolonization}
\subsection*{Unit 9: Globalization}

\section{Historical Thinking Skills and Reasoning Processes}
\begin{itemize}[leftmargin=2em]
    \item Developing historical arguments
    \item Analyzing sourcing and situation of documents
    \item Making comparisons across regions and time periods
    \item Explaining continuity and change over time
\end{itemize}

\section{Grading Breakdown}
\begin{center}
\begin{tabular}{L{0.55\linewidth}L{0.25\linewidth}}
\toprule
\textbf{Category} & \textbf{Percentage} \\
\midrule
AP-Style Tests and Essays (DBQ/LEQ) & 45\% \\
Quizzes and SAQs & 25\% \\
Projects and Presentations & 15\% \\
Homework and Participation & 15\% \\
\bottomrule
\end{tabular}
\end{center}

\chapter{United States History}

\noindent\textbf{Course:} United States History\\
\textbf{Grade Level:} 11th Grade\\
\textbf{Duration:} Full Year (36 weeks)\\
\textbf{Credits:} 1.0 Credit\\
\textbf{Prerequisites:} World History (recommended)

\section{Course Description}
United States History examines major political, economic, social, and cultural developments from
the colonial era to the present.

\section{Units of Study}
\subsection*{Unit 1: Colonization and Early America}
\subsection*{Unit 2: American Revolution and Constitution}
\subsection*{Unit 3: Expansion, Reform, and Civil War}
\subsection*{Unit 4: Reconstruction and Industrialization}
\subsection*{Unit 5: The United States as a World Power}
\subsection*{Unit 6: Depression, New Deal, and World War II}
\subsection*{Unit 7: Cold War and Civil Rights}
\subsection*{Unit 8: Contemporary America}

\section{Grading Breakdown}
\begin{center}
\begin{tabular}{L{0.55\linewidth}L{0.25\linewidth}}
\toprule
\textbf{Category} & \textbf{Percentage} \\
\midrule
Tests and Major Essays & 40\% \\
Quizzes & 20\% \\
Projects and Presentations & 20\% \\
Homework and Classwork & 15\% \\
Participation & 5\% \\
\bottomrule
\end{tabular}
\end{center}

\chapter{AP United States History}

\noindent\textbf{AP COURSE:} AP United States History\\
\textbf{Grade Level:} 11th--12th Grade\\
\textbf{Duration:} Full Year (36 weeks)\\
\textbf{Credits:} 1.0 Credit (Weighted 5.0 scale)\\
\textbf{Prerequisites:} Strong performance in prior social studies and English; teacher recommendation

\section{Course Description}
AP United States History (APUSH) is a college-level survey of U.S.\ history emphasizing the
development of historical arguments and the use of primary sources.

\section{AP Exam Components}
\begin{itemize}[leftmargin=2em]
    \item Multiple-choice questions on stimulus-based sets
    \item Short-answer questions
    \item DBQ and LEQ essays
\end{itemize}

\section{Periods and Themes}
The course is organized into nine historical periods (1491--present) and explores themes such as:
\begin{itemize}[leftmargin=2em]
    \item American and national identity
    \item Work, exchange, and technology
    \item Politics and power
    \item America in the world
\end{itemize}

\section{Grading Breakdown}
\begin{center}
\begin{tabular}{L{0.55\linewidth}L{0.25\linewidth}}
\toprule
\textbf{Category} & \textbf{Percentage} \\
\midrule
AP-Style Tests and Essays & 50\% \\
Quizzes and SAQs & 20\% \\
Projects and Presentations & 15\% \\
Homework and Participation & 15\% \\
\bottomrule
\end{tabular}
\end{center}

\chapter{US Government and Civics}

\noindent\textbf{Course:} US Government and Civics\\
\textbf{Grade Level:} 12th Grade\\
\textbf{Duration:} Semester (18 weeks) or Full Year (local decision)\\
\textbf{Credits:} 0.5--1.0 Credit\\
\textbf{Prerequisites:} United States History

\section{Course Description}
US Government and Civics introduces students to the structure and function of the American
political system, constitutional principles, and civic engagement.

\section{Course Goals}
\begin{itemize}[leftmargin=2em]
    \item Understand the foundations of American government
    \item Analyze the Constitution and landmark Supreme Court cases
    \item Explain the roles of the three branches of government
    \item Evaluate the impact of political parties, interest groups, and media
\end{itemize}

\section{Units of Study}
\subsection*{Unit 1: Foundations and Constitution}
\subsection*{Unit 2: Institutions of National Government}
\subsection*{Unit 3: Civil Rights and Civil Liberties}
\subsection*{Unit 4: Politics, Elections, and Public Policy}

\section{Grading Breakdown}
\begin{center}
\begin{tabular}{L{0.55\linewidth}L{0.25\linewidth}}
\toprule
\textbf{Category} & \textbf{Percentage} \\
\midrule
Tests and Essays & 40\% \\
Quizzes & 20\% \\
Projects and Presentations & 20\% \\
Homework and Classwork & 15\% \\
Participation & 5\% \\
\bottomrule
\end{tabular}
\end{center}

\chapter{AP United States Government and Politics}

\noindent\textbf{AP COURSE:} AP United States Government and Politics\\
\textbf{Grade Level:} 12th Grade\\
\textbf{Duration:} Full Year (36 weeks) or Semester (intensive)\\
\textbf{Credits:} 1.0 Credit (Weighted 5.0 scale)\\
\textbf{Prerequisites:} US History and strong reading/writing skills

\section{Course Description}
AP US Government and Politics is a college-level course on American political institutions,
behaviors, and policies, aligned with the College Board framework.

\section{AP Exam Components}
\begin{itemize}[leftmargin=2em]
    \item Multiple-choice questions
    \item Free-response questions including concept application, quantitative analysis, SCOTUS comparison, and argument essay
\end{itemize}

\section{Units of Study}
\subsection*{Unit 1: Foundations of American Democracy}
\subsection*{Unit 2: Interactions Among Branches of Government}
\subsection*{Unit 3: Civil Liberties and Civil Rights}
\subsection*{Unit 4: American Political Ideologies and Beliefs}
\subsection*{Unit 5: Political Participation}

\section{Required Documents and Supreme Court Cases}
Students study:
\begin{itemize}[leftmargin=2em]
    \item The US Constitution, Federalist and Anti-Federalist writings
    \item Required Supreme Court cases specified by the College Board
\end{itemize}

\section{Grading Breakdown}
\begin{center}
\begin{tabular}{L{0.55\linewidth}L{0.25\linewidth}}
\toprule
\textbf{Category} & \textbf{Percentage} \\
\midrule
AP-Style Tests and FRQs & 50\% \\
Quizzes and Reading Checks & 20\% \\
Projects and Presentations & 15\% \\
Homework and Participation & 15\% \\
\bottomrule
\end{tabular}
\end{center}

% =========================
% Part V: Additional Advanced Courses
% =========================
\part{Additional Advanced Courses}

\chapter{Dual Credit English Composition I \& II}

\noindent\textbf{DUAL CREDIT COURSE:} English Composition I (ENGL 1301) and English Composition II (ENGL 1302)\\
\textbf{Grade Level:} 11th--12th Grade\\
\textbf{Duration:} Semester each (or combined year-long)\\
\textbf{Credits:} High school English credit and college credit (through partner institution)\\
\textbf{Prerequisites:} College placement requirements and counselor/teacher approval

\section{Course Description}
Dual Credit English Composition I and II provide college-level instruction in academic writing,
reading, and research. Students earn both high school and college credit upon successful completion.

\section{Prerequisites and Placement}
Students must:
\begin{itemize}[leftmargin=2em]
    \item Meet college placement scores (TSI, SAT, ACT, or local standard)
    \item Complete application and registration with the partner college
\end{itemize}

\section{Composition II (ENGL 1302) Focus}
\begin{itemize}[leftmargin=2em]
    \item Advanced research and documentation
    \item Critical analysis of literature and non-fiction
    \item Longer, more complex argumentative essays
\end{itemize}

\chapter{Dual Credit Government and History}

\noindent\textbf{DUAL CREDIT COURSE:} US Government and US History\\
\textbf{Grade Level:} 11th--12th Grade\\
\textbf{Duration:} Semester courses (or year-long combination)\\
\textbf{Credits:} High school social studies credit and college credit\\
\textbf{Prerequisites:} College placement, counselor/teacher approval

\section{Course Description}
Dual Credit Government and History courses follow the curriculum of the partner college,
providing students with college-level experiences in US Government and US History while
fulfilling high school graduation requirements.

% =========================
% Appendices
% =========================
\appendix

\chapter{Course Planning Guide}

\section{Mathematics Pathways}
Sample pathways from Algebra 1 through AP Calculus and Statistics, including:
\begin{itemize}[leftmargin=2em]
    \item Standard pathway: Algebra 1 $\rightarrow$ Geometry $\rightarrow$ Algebra 2 $\rightarrow$ Pre-Calculus
    \item Advanced pathway: Honors Algebra 1 in 8th grade leading to AP Calculus BC
\end{itemize}

\section{Science Pathways}
Examples:
\begin{itemize}[leftmargin=2em]
    \item Biology $\rightarrow$ Chemistry $\rightarrow$ Physics
    \item Honors and AP sequences in Biology, Chemistry, and Physics
\end{itemize}

\section{Weighted GPA Comparison}
Explanation of GPA weighting:
\begin{itemize}[leftmargin=2em]
    \item Regular courses on a 4.0 scale
    \item Honors courses weighted (e.g., +0.5)
    \item AP and Dual Credit courses weighted (e.g., +1.0)
\end{itemize}

\section{College Admission Competitiveness}
Guidance on:
\begin{itemize}[leftmargin=2em]
    \item Course rigor and college admissions
    \item Recommended number of AP/dual credit courses for selective colleges
\end{itemize}

\section{Dual Credit vs.\ AP Comparison}
Overview of:
\begin{itemize}[leftmargin=2em]
    \item Credit transferability
    \item Cost differences
    \item Impact on college readiness and scheduling
\end{itemize}

\chapter{Summer Assignments}

\section{Common Summer Assignment Types}
Examples:
\begin{itemize}[leftmargin=2em]
    \item Required reading for AP English courses with written responses
    \item Pre-course problem sets for AP Math and Science
    \item Research or project proposals for advanced courses
\end{itemize}

\chapter{AP Exam Registration and Policies}

\section{Registration Timeline}
General guidance on:
\begin{itemize}[leftmargin=2em]
    \item Fall registration windows
    \item Late registration policies
\end{itemize}

\section{Exam Fees}
Information about:
\begin{itemize}[leftmargin=2em]
    \item Standard exam fees
    \item Fee reductions or waivers for eligible students
\end{itemize}

\section{Score Sending}
Overview of:
\begin{itemize}[leftmargin=2em]
    \item Designating colleges to receive scores
    \item Score reporting timelines
\end{itemize}

\section{College Credit Policies}
Notes on:
\begin{itemize}[leftmargin=2em]
    \item Variation in AP credit acceptance by institution
    \item Importance of checking college websites for current policies
\end{itemize}

\chapter*{Conclusion}
\addcontentsline{toc}{chapter}{Conclusion}

This comprehensive guide is designed to support students, families, and educators in planning
and pursuing rigorous, well-sequenced high school coursework aligned with college and career goals.

\end{document}
